\chapter*{Ek: Çoktan Seçmeli Sorular}
\addcontentsline{toc}{chapter}{Ek: Çoktan Seçmeli Sorular}
\markboth{Ek: Çoktan Seçmeli Sorular}{Ek: Çoktan Seçmeli Sorular}

Her bölüm için altı soru verilmiştir. Her soruda tek doğru seçenek vardır.

\newcommand{\coksoru}[5]{\item #1
  \begin{enumerate}
  \renewcommand{\labelenumii}{\textbf{\Alph{enumii})}}
    \item #2 \item #3 \item #4 \item #5
  \end{enumerate}}

\section*{Bölüm 1: İstatistiksel düşünme ve araştırma süreci}
\begin{enumerate}\renewcommand{\labelenumi}{B1.\arabic{enumi}.}
\coksoru{Bir evrene ait bilinmeyen sayısal özellik hangisidir?}{İstatistik}{Parametre}{Değişken}{Gözlem}
\coksoru{Her satırın bir öğrenciyi temsil ettiği veri dosyasında analiz birimi nedir?}{Ölçek maddesi}{Öğrenci}{Sütun}{Evren ortalaması}
\coksoru{Sayısal kod verilmesine rağmen nominal kalabilecek değişken hangisidir?}{Yaş}{Çalışma süresi}{Öğrenci numarası}{Test puanı}
\coksoru{Nedensel yorumu en güçlü biçimde destekleyen özellik hangisidir?}{Büyük korelasyon}{Düşük $p$-değeri}{Rastgele atama}{Büyük örneklem}
\coksoru{Bir yapının gözlenebilir değişkene dönüştürülmesine ne ad verilir?}{Tabakalama}{İşlemsel tanım}{Standartlaştırma}{Ağırlıklandırma}
\coksoru{Gönüllü katılımcılarla yapılan araştırmanın temel riski hangisidir?}{Yuvarlama hatası}{Seçim yanlılığı}{Çoklu bağlantı}{Sıra bağı}
\end{enumerate}

\section*{Bölüm 2: Örnekleme}
\begin{enumerate}\renewcommand{\labelenumi}{B2.\arabic{enumi}.}
\coksoru{Her birimin bilinen ve sıfırdan büyük seçilme olasılığı bulunduğu yöntem hangisidir?}{Kolayda örnekleme}{Olasılıklı örnekleme}{Kartopu örnekleme}{Amaçlı örnekleme}
\coksoru{Alt grupların ayrı temsil edilmesini sağlayan yöntem hangisidir?}{Tabakalı örnekleme}{Küme örneklemesi}{Sistematik hata}{Tam sayım}
\coksoru{Aynı sınıftaki öğrencilerin birbirine benzemesi en çok hangi sorunu doğurur?}{Bağımsız bilgi miktarının azalması}{Ortalamanın tanımsız olması}{Ölçeğin nominalleşmesi}{Medyanın değişmemesi}
\coksoru{Yanıtlamama yanlılığı ne zaman oluşur?}{Yanıt oranı yüzde 100 olduğunda}{Yanıtlamama sonuçla sistematik ilişkili olduğunda}{Örneklem büyük olduğunda}{Tabakalar eşit olduğunda}
\coksoru{Seçilme olasılığı düşük birimin analiz ağırlığı genellikle nasıldır?}{Daha büyük}{Sıfır}{Daima bir}{Daha küçük}
\coksoru{Örnekleme çerçevesinin hedef evrenden eksik olması hangi hatadır?}{Kapsama hatası}{Ölçüm değişmezliği}{Tip II hata}{Hesaplama hatası}
\end{enumerate}

\section*{Bölüm 3: Betimsel istatistik}
\begin{enumerate}\renewcommand{\labelenumi}{B3.\arabic{enumi}.}
\coksoru{Güçlü sağa çarpık dağılımda tipik konum için hangi ölçü daha sağlamdır?}{Ortalama}{Medyan}{Varyans}{Ranj}
\coksoru{Standart sapmanın karesi hangisidir?}{Varyans}{IQR}{Standart hata}{Değişim katsayısı}
\coksoru{$Q_1=10$ ve $Q_3=18$ ise IQR kaçtır?}{8}{14}{28}{4}
\coksoru{Bütün puanlara 5 eklendiğinde standart sapma nasıl değişir?}{5 artar}{Değişmez}{5 azalır}{İki katına çıkar}
\coksoru{Aykırı değerlerden en fazla etkilenen merkez ölçüsü hangisidir?}{Mod}{Medyan}{Ortalama}{Yüzdelik sıra}
\coksoru{Sınıflandırılmış veride yaklaşık ortalama hesaplanırken her sınıfı hangi değer temsil eder?}{Alt gerçek sınır}{Sınıf orta noktası}{Kümülatif frekans}{Sınıf genişliği}
\end{enumerate}

\section*{Bölüm 4: Normal dağılım ve standart puan}
\begin{enumerate}\renewcommand{\labelenumi}{B4.\arabic{enumi}.}
\coksoru{$z=1{,}5$ neyi gösterir?}{Puan ortalamanın 1,5 standart sapma üzerindedir}{Puan yüzde 1,5 artmıştır}{Olasılık 1,5'tir}{Puan medyana eşittir}
\coksoru{Standart normal dağılımın ortalama ve standart sapması nedir?}{0 ve 1}{1 ve 0}{50 ve 10}{0 ve 0}
\coksoru{Normal dağılımda yaklaşık yüzde 95'lik merkezî alan hangi sınırlar arasındadır?}{$-1$ ve $1$}{$-1{,}96$ ve $1{,}96$}{$0$ ve $1$}{$-3$ ve $3$}
\coksoru{Örneklem ortalamasının standart hatası hangisidir?}{$s\sqrt n$}{$s/n$}{$s/\sqrt n$}{$n/s$}
\coksoru{Örneklem hacmi dört katına çıkarsa standart hata yaklaşık ne olur?}{İki katına çıkar}{Yarıya iner}{Dört katına çıkar}{Değişmez}
\coksoru{Merkezî limit teoremi öncelikle hangi dağılımla ilgilidir?}{Ham puanların mutlaka normal olmasıyla}{Örneklem ortalamasının dağılımıyla}{Kategorilerin sırasıyla}{Aykırı değerlerin silinmesiyle}
\end{enumerate}

\section*{Bölüm 5: Nokta tahmini}
\begin{enumerate}\renewcommand{\labelenumi}{B5.\arabic{enumi}.}
\coksoru{Evren oranının doğal nokta tahmini hangisidir?}{$\bar x$}{$\hat p=x/n$}{$s^2/n$}{$z$}
\coksoru{Yansız tahmin edici için hangisi doğrudur?}{Her örneklemde parametreye eşittir}{Tekrarlı örneklemedeki beklenen değeri parametredir}{Varyansı sıfırdır}{Aykırı değerlerden etkilenmez}
\coksoru{MSE hangi iki bileşene ayrılır?}{Varyans ve yanlılığın karesi}{Ortalama ve medyan}{Ranj ve IQR}{Güç ve alfa}
\coksoru{Örneklem büyüdükçe hedef çevresinde yoğunlaşma hangi özelliktir?}{Tutarlılık}{Çarpıklık}{Kapsama}{Bağımsızlık}
\coksoru{Uç gözlemlerden az etkilenme hangi özelliktir?}{Sağlamlık}{Tamlık}{Yeterlilik}{Normallik}
\coksoru{Orantısız tabakalı örneklemde evren ortalaması için ne gerekir?}{Tabaka ağırlıkları}{Yalnız ham ortalama}{Kategori birleştirme}{Tek yönlü test}
\end{enumerate}

\section*{Bölüm 6: Güven aralıkları}
\begin{enumerate}\renewcommand{\labelenumi}{B6.\arabic{enumi}.}
\coksoru{Yüzde 95 güven düzeyinin doğru yorumu hangisidir?}{Parametrenin gerçekleşmiş aralıkta bulunma olasılığı yüzde 95'tir}{Yöntem tekrarlı örneklemede aralıkların yaklaşık yüzde 95'inde parametreyi kapsar}{Verilerin yüzde 95'i aralıktadır}{Hipotez yüzde 95 doğrudur}
\coksoru{Güven düzeyi artırıldığında aralık genellikle ne olur?}{Daralır}{Genişler}{Değişmez}{Sıfırlanır}
\coksoru{Standart hata azaldığında hata payı ne olur?}{Azalır}{Artar}{Daima değişmez}{Negatif olur}
\coksoru{Ortalama için küçük örneklemde $\sigma$ bilinmiyorsa hangi dağılım kullanılır?}{$t$ dağılımı}{Ki-kare uyum testi}{Binom dağılımı}{Fisher kesin testi}
\coksoru{Fark aralığının sıfırı içermesi ne gösterir?}{Grupların kesin eşit olduğunu}{Sıfır farkının veriyle uyumlu değerlerden biri olduğunu}{Etkinin kesin olmadığını}{Örneklemin yanlı olduğunu}
\coksoru{Güven aralığı hangi sorunu kapsamaz?}{Örnekleme değişkenliği}{Tahmin hassasiyeti}{Sistematik seçim yanlılığı}{Standart hata}
\end{enumerate}

\section*{Bölüm 7: Hipotez testleri ve etki}
\begin{enumerate}\renewcommand{\labelenumi}{B7.\arabic{enumi}.}
\coksoru{$p$-değeri hangisidir?}{$H_0$'ın doğru olma olasılığı}{$H_0$ altında gözlenen kadar uç sonuç olasılığı}{Etkinin büyüklüğü}{Testin gücü}
\coksoru{Doğru $H_0$'ı reddetmek hangi hatadır?}{Tip I}{Tip II}{Örnekleme çerçevesi}{Regresyon hatası}
\coksoru{Güç hangi ifadeye eşittir?}{$1-\alpha$}{$1-\beta$}{$\alpha+\beta$}{$p/2$}
\coksoru{Tek yönlü test ne zaman seçilmelidir?}{Sonuç görüldükten sonra}{Yön önceden gerekçelendirildiğinde}{Her zaman}{İki yön de önemli olduğunda}
\coksoru{Büyük örneklemde çok küçük etki neden anlamlı çıkabilir?}{Standart hata küçülebildiği için}{Etki büyüdüğü için}{Yanlılık yok olduğu için}{Medyan değiştiği için}
\coksoru{Çoklu testin temel riski hangisidir?}{Yanlış pozitiflerin artması}{Örneklemin küçülmesi}{Ortalamanın tanımsızlaşması}{Varyansın sıfırlanması}
\end{enumerate}

\section*{Bölüm 8: t testleri}
\begin{enumerate}\renewcommand{\labelenumi}{B8.\arabic{enumi}.}
\coksoru{Aynı öğrencilerin ön-test ve son-test puanları için uygun test hangisidir?}{Bağımsız $t$}{Eşleştirilmiş $t$}{Tek yönlü ANOVA}{Ki-kare}
\coksoru{Bir örneklem ortalamasını norm değeriyle karşılaştıran test hangisidir?}{Tek örneklem $t$}{Welch ANOVA}{Mann--Whitney}{Pearson korelasyonu}
\coksoru{Welch testi hangi varsayımı gerektirmez?}{Bağımsız gözlemler}{Nicel sonuç}{Eşit evren varyansları}{Uygun tasarım}
\coksoru{Eşleştirilmiş testin temel analiz değeri hangisidir?}{İki ayrı grup ortalaması}{Bireysel fark puanı}{Kategori frekansı}{Sıra toplamı}
\coksoru{$t$ istatistiğinin paydasında ne bulunur?}{Tahminin standart hatası}{Örneklem ortalaması}{Etki büyüklüğü}{Anlamlılık düzeyi}
\coksoru{Cohen'in $d$ değeri neyi özetler?}{Standartlaştırılmış ortalama farkını}{Tip I hata olasılığını}{Örneklem oranını}{Serbestlik derecesini}
\end{enumerate}

\section*{Bölüm 9: ANOVA}
\begin{enumerate}\renewcommand{\labelenumi}{B9.\arabic{enumi}.}
\coksoru{Genel ANOVA sıfır hipotezi hangisidir?}{Bütün ortalamalar eşittir}{Bütün ortalamalar farklıdır}{Varyanslar sıfırdır}{Medyanlar aynıdır}
\coksoru{$F$ istatistiği hangi orandır?}{Gruplar arası ortalama kare / grup içi ortalama kare}{Toplam / örneklem hacmi}{Medyan / IQR}{Korelasyon / eğim}
\coksoru{Anlamlı genel $F$ testi neyi göstermez?}{En az bir fark için kanıtı}{Hangi grupların farklı olduğunu}{Ortak hipotezin uyumsuzluğunu}{Gruplar arası değişkenliği}
\coksoru{Benzer varyanslarla bütün ikililer için hangi yöntem uygundur?}{Tukey HSD}{Games--Howell}{İşaret testi}{Fisher dönüşümü}
\coksoru{Eşit olmayan varyanslarda hangi ikili yöntem uygundur?}{Games--Howell}{Tukey'in zorunlu kullanımı}{Pearson ki-kare}{Tek örneklem $t$}
\coksoru{Eta-kare neyi özetler?}{Toplam değişkenliğin gruplarla ilişkili oranını}{Yanlış pozitif oranını}{Bireysel tahmin hatasını}{Medyan farkını}
\end{enumerate}

\section*{Bölüm 10: Korelasyon}
\begin{enumerate}\renewcommand{\labelenumi}{B10.\arabic{enumi}.}
\coksoru{Pearson korelasyonu öncelikle hangi ilişkiyi özetler?}{Doğrusal ilişkiyi}{Her tür nedenselliği}{Yalnız medyan farkını}{Kategori bağımsızlığını}
\coksoru{$r=0$ ne anlama gelir?}{Hiçbir ilişki yoktur}{Doğrusal ilişki yoktur; eğrisel ilişki olabilir}{Değişkenler bağımsızdır}{Etki sıfırdır}
\coksoru{Spearman korelasyonu hangi bilgiye dayanır?}{Sıralara}{Yalnız ortalamalara}{Hücre frekanslarına}{Varyans oranına}
\coksoru{$r=-0{,}70$ için $r^2$ kaçtır?}{0,49}{-0,49}{0,70}{1,40}
\coksoru{Korelasyonun nedensellik göstermemesinin nedeni hangisidir?}{Ters nedensellik ve karıştırıcılar bulunabilmesi}{Katsayının birimsiz olması}{Değerin negatif olabilmesi}{Örneklem hacminin yazılması}
\coksoru{Korelasyon analizinin ilk grafiği hangisidir?}{Saçılım grafiği}{Pasta grafiği}{Yalnız çubuk grafik}{Akış şeması}
\end{enumerate}

\section*{Bölüm 11: Basit doğrusal regresyon}
\begin{enumerate}\renewcommand{\labelenumi}{B11.\arabic{enumi}.}
\coksoru{Regresyon eğimi neyi gösterir?}{X bir birim arttığında beklenen Y ortalamasındaki değişimi}{Y'nin örneklem hacmini}{X'in medyanını}{Tip I hatayı}
\coksoru{Artık nasıl hesaplanır?}{Gözlenen Y eksi uydurulan Y}{Uydurulan Y eksi X}{Ortalama eksi medyan}{SST artı SSE}
\coksoru{En küçük kareler hangi toplamı küçültür?}{Artık kareleri toplamını}{Mutlak X toplamını}{Korelasyonların toplamını}{Örneklem hacmini}
\coksoru{Basit regresyonda $R^2$ ile korelasyon arasındaki ilişki hangisidir?}{$R^2=r^2$}{$R^2=r$}{$R^2=1-r$}{$R^2=r/2$}
\coksoru{Bireysel tahmin aralığı neden ortalama güven aralığından geniştir?}{Bireysel artık değişkenliğini de içerir}{Daha düşük güven kullanır}{Örneklemi küçültür}{Eğimi sıfırlar}
\coksoru{Gözlenen X aralığının dışında tahmin yapmaya ne denir?}{Ekstrapolasyon}{Merkezleme}{Tabakalama}{Standartlaştırma}
\end{enumerate}

\section*{Bölüm 12: Kategorik veriler ve ki-kare}
\begin{enumerate}\renewcommand{\labelenumi}{B12.\arabic{enumi}.}
\coksoru{Bağımsızlık altında beklenen hücre frekansı nasıl bulunur?}{Satır toplamı çarpı sütun toplamı bölü genel toplam}{İki yüzdeyi toplama}{Yalnız satır toplamı}{Gözlenen frekansın karesi}
\coksoru{$r\times c$ tabloda serbestlik derecesi hangisidir?}{$(r-1)(c-1)$}{$rc$}{$r+c$}{$N-1$}
\coksoru{Küçük $2\times2$ tabloda hangi yöntem düşünülebilir?}{Fisher kesin testi}{Pearson korelasyonu}{Welch ANOVA}{Eşleştirilmiş $t$}
\coksoru{Cramér $V$ neyi özetler?}{Kategorik ilişkinin büyüklüğünü}{Ortalama farkını}{Regresyon eğimini}{Standart hatayı}
\coksoru{Grup büyüklükleri farklıysa tabloyu yorumlamak için ne gerekir?}{Uygun satır veya sütun yüzdeleri}{Yalnız ham sayılar}{Yalnız genel toplam}{Normal dağılım eğrisi}
\coksoru{Hangi hücrelerin genel sonuca katkı verdiğini ne gösterir?}{Standartlaştırılmış artıklar}{Yalnız $p$-değeri}{Örneklem ortalaması}{Regresyon kesmesi}
\end{enumerate}

\section*{Bölüm 13: Parametrik olmayan yöntemler}
\begin{enumerate}\renewcommand{\labelenumi}{B13.\arabic{enumi}.}
\coksoru{İki bağımsız sıralı grup için hangi test uygundur?}{Mann--Whitney}{Wilcoxon işaretli sıra}{Eşleştirilmiş $t$}{McNemar}
\coksoru{Aynı kişilerin iki ölçümünde sıra büyüklüklerini kullanan test hangisidir?}{Wilcoxon işaretli sıra}{Mann--Whitney}{Kruskal--Wallis}{Pearson ki-kare}
\coksoru{Üç bağımsız sıralı grup için hangi test uygundur?}{Kruskal--Wallis}{İşaret testi}{Tek örneklem $t$}{Pearson korelasyonu}
\coksoru{Mann--Whitney her zaman neyin testi değildir?}{Medyan eşitliğinin}{Dağılım farkının}{Sıra üstünlüğünün}{Bağımsız grupların}
\coksoru{İşaret testi hangi bilgiyi kullanır?}{Farkların yalnız yönünü}{Farkların karelerini}{Grup varyanslarını}{Korelasyon eğimini}
\coksoru{Parametrik olmayan testler için hangisi doğrudur?}{Bağımsızlık gibi varsayımları vardır}{Hiç varsayımları yoktur}{Her zaman daha güçlüdür}{Her zaman ortalama farkını sınar}
\end{enumerate}

\section*{Bölüm 14: Bütünleştirici laboratuvar}
\begin{enumerate}\renewcommand{\labelenumi}{B14.\arabic{enumi}.}
\coksoru{Analiz akışında ilk adım hangisidir?}{Araştırma sorusunu tanımlamak}{$p$-değerini hesaplamak}{Sonucu yazmak}{Aykırı değer silmek}
\coksoru{Ham veri üzerinde elle ve kayıtsız değişiklik yapmak neyi bozar?}{Yeniden üretilebilirliği}{Örneklem ortalamasının tanımını}{Kategori sayısını}{Normal dağılımı}
\coksoru{Aynı öğrencinin ön ve son puanının aynı satırda olması neyi korur?}{Eşleştirme bilgisini}{Bağımsız grup yapısını}{Rastgele atamayı}{Evren büyüklüğünü}
\coksoru{Çok yüksek korelasyon görüldüğünde ilk yapılması gereken hangisidir?}{Grafik, kodlama ve etkili gözlemleri denetlemek}{Nedensellik ilan etmek}{Veriyi kategorileştirmek}{Örneklemi küçültmek}
\coksoru{SPSS analizinin yeniden üretilebilirliğini ne artırır?}{Syntax dosyasını saklamak}{Yalnız Viewer çıktısını saklamak}{Menüleri hatırlamak}{Ekran görüntüsü almak}
\coksoru{Birincil analiz sonucu ne zaman belirlenmelidir?}{Veri görülmeden önce}{En küçük $p$ bulunduktan sonra}{Rapor yazılırken}{Hakem istedikten sonra}
\end{enumerate}